\chapter{Аналитическая часть}

\section{Описание системы}

Разрабатываемая система представляет собой веб-портал для бронирования автомобилей на выбранные пользователем даты. Пользователь может зарегистрироваться, пройти авторизацию, просмотреть список автомобилей, выбрать период аренды, оформить бронирование и оплатить его. После оформления пользователь может просматривать свои бронирования, завершать активную аренду или отменять ее.

Система построена как набор микросервисов. Каждый микросервис отвечает за отдельную задачу и имеет доступ только к собственной области данных. Взаимодействие пользователя с бизнес-функциями выполняется через сервис-агрегатор, который координирует вызовы сервисов автомобилей, бронирования, платежей, статистики и авторизации.

Для авторизации используется собственный Identity Provider, реализующий OpenID Connect. Пользовательский интерфейс рассматривается как стороннее приложение, получающее токен через Authorization Code Flow. Все защищенные запросы выполняются с JWT-токеном, который валидируется сервисами по JWKs.

\section{Функциональные требования}

Система должна обеспечивать следующие функции:

\begin{enumerate}
  \item Регистрация новых пользователей через пользовательский интерфейс.
  \item Авторизация пользователей через собственный Identity Provider.
  \item Разделение пользователей на роли:
  \begin{itemize}
    \item гость;
    \item пользователь;
    \item администратор.
  \end{itemize}
  \item Просмотр списка автомобилей.
  \item Фильтрация автомобилей по марке, модели, номеру и типу.
  \item Выбор дат аренды автомобиля.
  \item Создание бронирования автомобиля на выбранные даты.
  \item Создание платежа при бронировании автомобиля.
  \item Просмотр списка собственных бронирований.
  \item Завершение активного бронирования.
  \item Отмена бронирования.
  \item Создание пользователей администратором.
  \item Просмотр администратором списка пользователей.
  \item Просмотр администратором статистики, собранной по событиям системы.
\end{enumerate}

Функции пользователей представлены в таблице \ref{tbl:user-func}.

\begin{longtable}{|p{0.5cm}|p{15.5cm}|}
	\caption{Функции пользователей}
	\label{tbl:user-func} \\
	\hline
	\begin{rotatebox}[origin=r]{90}{\textbf{Гость}}\end{rotatebox}
	&
	1. Просмотр страницы входа; \newline
	2. Регистрация в системе; \newline
	3. Переход к форме авторизации Identity Provider. \\
	\hline
	\begin{rotatebox}[origin=r]{90}{\textbf{Пользователь}}\end{rotatebox}
	&
	1. Авторизация в системе; \newline
	2. Просмотр списка автомобилей; \newline
	3. Поиск автомобилей по основным полям; \newline
	4. Выбор дат аренды; \newline
	5. Бронирование автомобиля; \newline
	6. Автоматическое создание платежа; \newline
	7. Просмотр собственных бронирований; \newline
	8. Завершение активной аренды; \newline
	9. Отмена бронирования. \\
	\hline
	\begin{rotatebox}[origin=r]{90}{\textbf{Администратор}}\end{rotatebox}
	&
	1. Все функции пользователя; \newline
	2. Создание новых пользователей; \newline
	3. Просмотр списка пользователей; \newline
	4. Просмотр отчета статистики; \newline
	5. Просмотр последних событий системы. \\
	\hline
\end{longtable}

\section{Входные данные}

Входные параметры системы представлены в таблице \ref{tbl:input}.

\begin{longtable}{|p{4cm}|p{12cm}|}
	\caption{Входные данные}
	\label{tbl:input} \\
	\hline
	\textbf{Операция} & \textbf{Входные данные} \\
	\hline
	\endfirsthead
	\hline
	\textbf{Операция} & \textbf{Входные данные} \\
	\hline
	\endhead
	\hline
	\multicolumn{2}{c}{\textit{Продолжение на следующей странице}}
	\endfoot
	\hline
	\endlastfoot

	Регистрация пользователя
	&
	1. \textit{логин}; \newline
	2. \textit{email}; \newline
	3. \textit{имя пользователя}; \newline
	4. \textit{пароль}. \\
	\hline

	Авторизация пользователя
	&
	1. \textit{логин}; \newline
	2. \textit{пароль}; \newline
	3. \textit{client\_id}; \newline
	4. \textit{redirect\_uri}; \newline
	5. \textit{scope}: openid, profile, email. \\
	\hline

	Получение списка автомобилей
	&
	1. \textit{номер страницы}; \newline
	2. \textit{размер страницы}; \newline
	3. \textit{признак отображения занятых автомобилей}; \newline
	4. \textit{поисковая строка} в пользовательском интерфейсе. \\
	\hline

	Создание бронирования
	&
	1. \textit{идентификатор автомобиля}; \newline
	2. \textit{дата начала аренды}; \newline
	3. \textit{дата окончания аренды}; \newline
	4. \textit{JWT-токен пользователя}. \\
	\hline

	Завершение или отмена бронирования
	&
	1. \textit{идентификатор бронирования}; \newline
	2. \textit{JWT-токен пользователя}. \\
	\hline

	Создание пользователя администратором
	&
	1. \textit{логин}; \newline
	2. \textit{email}; \newline
	3. \textit{имя пользователя}; \newline
	4. \textit{пароль}; \newline
	5. \textit{JWT-токен администратора}. \\
	\hline
\end{longtable}

\section{Выходные данные}

Выходными данными системы являются веб-страницы и JSON-ответы API. Основные выходные данные представлены в таблице \ref{tbl:output-data}.

\begin{longtable}{|p{0.5cm}|p{15.5cm}|}
	\caption{Выходные данные}
	\label{tbl:output-data} \\
	\hline
	\begin{rotatebox}[origin=r]{90}{\textbf{Гость}}\end{rotatebox}
	&
	1. Страница входа; \newline
	2. Форма регистрации; \newline
	3. Форма авторизации Identity Provider. \\
	\hline
	\begin{rotatebox}[origin=r]{90}{\textbf{Пользователь}}\end{rotatebox}
	&
	1. Список автомобилей: марка, модель, регистрационный номер, мощность, тип, цена за день, доступность; \newline
	2. Расчет стоимости аренды за выбранный период; \newline
	3. Список бронирований: автомобиль, даты, статус бронирования, статус и сумма платежа; \newline
	4. Сообщения об успешном бронировании, завершении или отмене аренды. \\
	\hline
	\begin{rotatebox}[origin=r]{90}{\textbf{Администратор}}\end{rotatebox}
	&
	1. Список пользователей: логин, email, имя, роль; \newline
	2. Отчет статистики: количество событий, количество типов действий, количество пользователей в отчете; \newline
	3. Последние события системы: время, пользователь, действие, данные события. \\
	\hline
\end{longtable}

\section{Состав системы}

Система состоит из следующих сервисов:

\begin{itemize}
  \item \textbf{ui-service} --- сервис пользовательского интерфейса, реализующий Single Page Application на React;
  \item \textbf{identity-service} --- сервис авторизации и данных пользовательских аккаунтов, выполняющий роль Identity Provider;
  \item \textbf{cars-service} --- сервис автомобилей;
  \item \textbf{rental-service} --- сервис бронирования автомобилей;
  \item \textbf{payment-service} --- сервис платежей;
  \item \textbf{statistics-service} --- сервис сбора и предоставления статистики;
  \item \textbf{gateway-service} --- сервис агрегирования запросов и предоставления ограниченного API для внешних клиентов;
  \item \textbf{kafka} --- брокер сообщений Apache Kafka для асинхронной передачи событий.
\end{itemize}

\section{Требования к безопасности}

Все методы \texttt{/api/**}, кроме публичных методов регистрации, авторизации и callback, должны быть закрыты token-based авторизацией. В качестве токена используется JWT, полученный от Identity Provider. Каждый сервис валидирует токен самостоятельно по JWKs. Пользователь определяется по claims токена, заголовок \texttt{X-User-Name} для определения пользователя не используется.

При отсутствии токена, ошибке валидации JWT или истечении срока действия токена сервис должен вернуть HTTP-статус 401.

\section{Требования к отказоустойчивости и наблюдаемости}

Система должна обрабатывать ошибки зависимых сервисов. Некритичный функционал должен деградировать без полного отказа пользовательского сценария. Например, при временной недоступности сервиса статистики основная операция бронирования не должна завершаться ошибкой.

Все сервисы должны логгировать входящие запросы и основные бизнес-события. Действия пользователя дополнительно отправляются через Kafka в сервис статистики.

\section{Вывод}

В аналитической части были определены назначение системы, роли пользователей, функциональные требования, входные и выходные данные, состав микросервисов, а также требования к безопасности и отказоустойчивости.
